% TOP_PROBLEM: {"id": 3096, "title": "The Hodge Conjecture", "category_id": 6, "category": {"id": 6, "name": "geometry", "display_name": "Geometry", "description": "Euclidean and non-Euclidean geometry, geometric structures.", "slug": "geometry", "order_index": 6, "created_at": "2026-07-31T15:26:25.671Z"}, "status": "open", "problem_number": "OPG-1803", "set_id": 12, "statusQualification": "The standard rational Hodge conjecture is used, explicitly restoring smoothness of the ambient projective variety."}
% FIRST_POSTED: 2026-10-01
% ENTRY_KIND: statement_only
% UnsolvedMath ID 3096; source catalogue code OPG-1803
% !TeX program = lualatex
% Definition and sources only; this file is not an LLM solution attempt.
\documentclass[11pt,a4paper]{article}
\usepackage[margin=25mm]{geometry}
\usepackage{fontspec}
\setmainfont{lmroman10-regular.otf}[BoldFont=lmroman10-bold.otf,
 ItalicFont=lmroman10-italic.otf,BoldItalicFont=lmroman10-bolditalic.otf]
\usepackage{amsmath,amssymb,mathtools}
\usepackage{unicode-math}
\setmathfont{latinmodern-math.otf}
\AtBeginDocument{\let\setminus\smallsetminus}
\usepackage{xurl,hyperref}
\hypersetup{colorlinks=true,urlcolor=blue}
\setcounter{secnumdepth}{0}
\setlength{\parindent}{0pt}
\setlength{\parskip}{5pt}
\setlength{\emergencystretch}{3em}
\begin{document}
\section*{The Hodge Conjecture}\label{3096}
{\small UnsolvedMath ID 3096 · OPG-1803 · Geometry · OpenGarden}
\subsection{Definitions and mathematical statement}
Let $X$ be a smooth projective algebraic variety over $\mathbb C$ of
complex dimension $d$. Thus $X$ is a nonsingular closed algebraic subvariety
of some complex projective space. For $0\le p\le d$, its singular
cohomology with complex coefficients has Hodge decomposition
\[
 H^{2p}(X,\mathbb C)=\bigoplus_{a+b=2p}H^{a,b}(X).
\]
The summand $H^{a,b}(X)$ is represented by harmonic differential forms
of complex bidegree $(a,b)$, using a Kähler metric.
Via extension of coefficients, view $H^{2p}(X,\mathbb Q)$ as a rational
subspace of $H^{2p}(X,\mathbb C)$. A rational Hodge class of codimension
$p$ is an element
$\alpha\in H^{2p}(X,\mathbb Q)\cap H^{p,p}(X)$.

An irreducible closed algebraic subvariety $Z\subseteq X$ of complex
codimension $p$ has a fundamental cycle, whose Poincaré dual gives a
cohomology class $[Z]\in H^{2p}(X,\mathbb Q)$. The subvariety $Z$ itself
need not be smooth.
The rational Hodge conjecture asserts that every rational Hodge class
can be expressed as
\[
 \alpha=\sum_{j=1}^m q_j[Z_j],\qquad q_j\in\mathbb Q,
\]
for a finite list of codimension-$p$ algebraic subvarieties of $X$.
All dimensions, all $p$ in the indicated range, and all such classes
are quantified over.

Smoothness of $X$ is an essential scope condition omitted from the short
catalogue statement and supplied here in the classical formulation.
Rational coefficients are essential too: the assertion neither demands
a single effective subvariety nor requires integral coefficients in the
linear combination.
\subsection{Short English statement}
On a smooth complex projective variety, is every rational cohomology class of type $(p,p)$ a rational combination of algebraic subvariety classes?
\subsection{Sources}
\begin{itemize}
\item[S1] ULAM.AI and the UnsolvedMath Contributors, supplied problems.json, record 3096 (OPG-1803), OpenGarden collection. Catalogue identity and source transcription; CC BY 4.0.
\url{https://huggingface.co/datasets/ulamai/UnsolvedMath}.
\item[S2] Open Problem Garden, The Hodge Conjecture. Original problem and discussion; the mathematical formulation below is expanded from the supplied source record.
\url{http://www.openproblemgarden.org/op/the_hodge_conjecture}.
\item[S3] P. Deligne, The Hodge Conjecture, official problem description, Clay Mathematics Institute.
\url{https://www.claymath.org/wp-content/uploads/2022/06/hodge.pdf}.
\end{itemize}
\end{document}
